problem:  
Let \( X \) be a compact Calabi–Yau threefold equipped with a Ricci-flat Kähler metric \( g \) and holomorphic volume form \( \Omega \).  
Suppose \( f_0 : X \to B_0 \) is a special Lagrangian torus fibration with singular fibers confined to a discriminant locus \( \Delta_0 \subset B_0 \), and assume that \( f_0 \) is **semi-flat** (i.e., smooth) away from \( \Delta_0 \).  
Let \( \{g_t\}_{t \in (-\epsilon, \epsilon)} \) be a smooth 1-parameter family of Ricci-flat Kähler metrics on \( X \) lying in a fixed Kähler class, and \( \{\Omega_t\} \) the corresponding holomorphic volume forms.  

Is it true that there exists a family of maps  
\[
f_t : X \to B_t 
\]
defined for sufficiently small \( |t| \), such that  
1. each \( f_t \) is a special Lagrangian torus fibration over \( B_t \setminus \Delta_t \), smooth away from a controlled singular set \( \Delta_t \),  
2. \( f_t \to f_0 \) smoothly (in the semi-flat region) as \( t \to 0 \), and  
3. the discriminant loci \( \Delta_t \) vary continuously (in a suitable geometric sense) with \( t \)?  

In other words, can small Ricci-flat deformations of a Calabi–Yau preserve (up to controlled singularities) the existence of a nearby special Lagrangian fibration structure?  

Why is it a "good" problem:  
This formulation integrates several central ideas: Ricci-flat Kähler geometry, deformation theory, and calibrated geometry, while integrating them in a **realistic** and **open** setting deeply tied to the **geometric formulation of mirror symmetry** (the Strominger–Yau–Zaslow conjecture). Unlike unrealistic global-foliation statements, this problem narrows focus to the stability of existing special Lagrangian fibrations under small perturbations of the Ricci-flat metric—a setting where known tools (e.g., McLean’s deformation theory, geometric analysis of Monge–Ampère equations, and metric collapse techniques) might apply but do not yet give a complete answer.  
The problem straddles differential geometry, symplectic topology, and complex deformation theory; resolving it would clarify how mirror symmetry behaves under small geometric deformations. It is simple to state yet profoundly deep, illustrating the essence of a “good” research-level question: conceptually crisp, bridging distinct fields, and guiding progress toward an open frontier.